\documentclass[a4paper]{article}
\usepackage[margin=2.2cm]{geometry}
\usepackage{parskip}
\title{Direct-manipulation TikZ iPad editor}
\date{2026-09-13}
\author{}
\begin{document}
\maketitle

\section*{Change}
The iPad TikZ editor now behaves as a drawing canvas instead of a form with a permanent right-side inspector. Rectangle and ellipse tools are selected from the compact toolbar and the user drags directly on the canvas to determine position and size. Connectors remain node-to-node and are rendered interactively on the canvas.

Element properties are shown directly below the drawing-tool row. Node fill color, border color, and border width are editable there; connector label, line color, line width, and arrow state are editable in the same location. Supported edits continue to rewrite only the conservative named-node/named-edge TikZ subset and preserve unrelated options.

The header is compact: it shows only the containing folder plus filename, then Render, a gear menu containing Auto layout and Source, and Close. The former right inspector and duplicated Auto layout/Source/Render toolbar buttons were removed.

\section*{Safety}
The canonical TikZ source remains the file on disk. The live SwiftUI canvas is a direct visual projection of that source, while the Ubuntu LuaLaTeX render remains available from the play button as a verification render rather than blocking basic drawing interaction.

\end{document}
